% Nouveau EGA V (cote 25): archive cote 25, site page 56 (PDF page 57).
% typescript. Draft compared: none.
% AI-transcribed by Claude Opus 5.5 (claude-opus-5-5, adaptive thinking, effort high), Message Batches API, batch ega25b (msgbatch_01A7n5YuXubUPNK1zsepi8RG), 2026-11-13; checked against both supplied scan parts; see CHECK_DECISIONS.tsv.
\scanpage{25/56}
\mspage{51}

n  au lieu de \xout{secti\unsure{o|n}} hyperplans.

\uline{Proposition} 9.3. \xout{et} Supposons \inserted[typed]{X géom irréd et} dim f(X) $\hand{\geq}$ \unsure{2}, \xout{et} soit x$\hand{\in}$X(k), \inserted{Soit $n\geq 2$} \xout{et}
\xout{\unsure{soit}} \inserted[typed]{a\hand{)}} \deleted{Alors pour n$\hand{\geq}$2}, \over{c}{C}onsidér\over{a}{o}n\over{t}{s} la famille \inserted{(linéaire)} des sections hyperplanes
de X relativement à $f_n=u_nf$ \uline{qui passent par} x, son élément
générique \xout{est} définit une $Y^{(n)}_{\hand{\unsure{\xi|\zeta}}}$ qui est \uline{géométriquement} \xout{intègre}irréd.

b) \xout{Si f(X) n'est pas contenu dans un plan défini sur $\overset{\hand{\scriptscriptstyle -}}{k}$, et n$\hand{\geq}$2,}
\xout{ou si n$\hand{\geq}$3,} Soit L une sous-variété linéaire de P passant par f(x) et
ne contenant pas f(X). \deleted{Alors pour n$\hand{\geq}$2}, \over{c}{C}onsidér\over{a}{o}n\over{t}{s} la famille \inserted[typed]{(linéaire)} des
sections hyperplanes de\xout{s} X relativement à \xout{\unsure{l'immer}} $f_n$ \xout{qui pa} , définis
par des hyperplans de P(n) \xout{contenant} \uline{``tangents à L en x''} (i.e. définis
par des \inserted{n-}formes sur P qui \inserted[typed]{sur L} s'annulent à l'ordre 2 au moins en x), son
membre générique \deleted{définit} \inserted{est} une $Y^{(n)}_{\hand{\unsure{\xi|\zeta}}}$ géométriquement irréductible.

c) Supposons X lisse en x, \xout{et f non} et \deleted{soit} n$\hand{\geq}$3, \inserted{ou f(X) pas contenu dans un plan défini sur $\bar{k}$.} \xout{alors} \over{c}{C}onsidérons
la famille des sections hyperplanes \inserted{\deleted{\unsure{tangentielles}}} \xout{génériques} $Y^{(n)}_{\hand{\unsure{\xi|\zeta}}}$ de X relativement
à $f_n$ qui sont \hand{``}\uline{tangentes à X en x}\hand{''}\unsure{\over{.}{,}} Alors le membre générique de ladite
définit une $Y^{(n)}_{\hand{\unsure{\xi|\zeta}}}$ géométriquement irréductible.

La démonstration est essentiellement triviale en termes des
critères de la fin du n\textsuperscript{o} précédent. Prenant un modèle affine de P
contenant f(x), on est ramené dans a) à trouver \xout{une} trois polynômes
en les coordonnées $\unsure{T|X}_1,..,\unsure{T|X}_r$ , de degré \over{\illegible{1}}{$\leq$}2, soient P,Q,R, tels que
Q(t)/P(t) et R(t)/P(t) soient alg indép sur k dans K (où K est le
corps des fonctions de X, et t = $(t_1,..,t_r)$ est le système d'élé-%
ments de K défini par les $T_i$ ); dans b), on exige de plus que P,Q,R
s'annulent à l'ordre 2 au moins sur L, qu'on peut d'ailleurs supposer
défini par des équations $T_1,..T_s = 0$ ; enfin dans \unsure{c|b}) c'est le même
pb, mais L étant l'image de l'espace tangent à X en x, et on se permet
\inserted{éventuellement} de prendre P,Q,R \xout{seulement} de degré \hand{3} , i.e. un peu plus de marge.

\mspage{56}
